\textbf{Phần 1. Trắc nghiệm lựa chọn}
\setcounterpageref{ex}{0}

\begin{ex} %[3D3V34]
Ước lượng chiều dài của một chiếc bút chì gỗ tiêu chuẩn là khoảng
\choice
{$18\text{ mm}$}
{$18\text{ m}$}
{$18\text{ dm}$}
{\True $18\text{ cm}$}
\loigiai{
Bút chì gỗ thông thường dài khoảng $15$ đến $18\text{ cm}$.
\textbf{Đáp án của câu hỏi là D}
}
\end{ex}

\begin{ex} %[3D3V34]
Ước lượng khối lượng của một quả trứng gà là khoảng
\choice
{$500\text{ g}$}
{\True $50\text{ g}$}
{$5\text{ kg}$}
{$5\text{ g}$}
\loigiai{
Một quả trứng gà nặng khoảng $50$ đến $60\text{ g}$.
\textbf{Đáp án của câu hỏi là B}
}
\end{ex}

\begin{ex} %[3D3V34]
Ước lượng dung tích của một thìa cà phê nhỏ là khoảng
\choice
{$50\text{ ml}$}
{$500\text{ ml}$}
{\True $5\text{ ml}$}
{$1\ell$}
\loigiai{
Một thìa cà phê chứa khoảng $5\text{ ml}$ chất lỏng.
\textbf{Đáp án của câu hỏi là C}
}
\end{ex}

\begin{ex} %[3D3V34]
Khi đo độ dày của một chiếc thẻ ngân hàng (hoặc thẻ học sinh), đơn vị đo thích hợp nhất là
\choice
{\True Mi-li-mét ($\text{mm}$)}
{Xăng-ti-mét ($\text{cm}$)}
{Mét ($\text{m}$)}
{Đề-xi-mét ($\text{dm}$)}
\loigiai{
Độ dày của thẻ rất mỏng (khoảng $1\text{ mm}$) nên dùng mi-li-mét là chuẩn xác nhất.
\textbf{Đáp án của câu hỏi là A}
}
\end{ex}

\begin{ex} %[3D3V34]
Một bình sữa đầy $250\text{ ml}$. Bé uống một nửa bình sữa. Bé đã uống bao nhiêu mi-li-lít sữa?
\choice
{$100\text{ ml}$}
{$150\text{ ml}$}
{$200\text{ ml}$}
{\True $125\text{ ml}$}
\loigiai{
Bé uống một nửa: $250 : 2 = 125\text{ (ml)}$.
\textbf{Đáp án của câu hỏi là D}
}
\end{ex}

\begin{ex} %[3D3V34]
Bác thợ mộc đo chiều dài mép bàn được $1\text{ m } 20\text{ cm}$. Chiều dài đó bằng bao nhiêu mi-li-mét?
\choice
{$120\text{ mm}$}
{\True $1200\text{ mm}$}
{$12000\text{ mm}$}
{$1020\text{ mm}$}
\loigiai{
Đổi $1\text{ m } 20\text{ cm} = 120\text{ cm} = 1200\text{ mm}$.
\textbf{Đáp án của câu hỏi là B}
}
\end{ex}

\begin{ex} %[3D3V34]
Mẹ mua một túi táo cân nặng $1\text{ kg } 200\text{ g}$. Mẹ lấy ra 2 quả táo cân được $350\text{ g}$. Số táo còn lại trong túi nặng là
\choice
{$950\text{ g}$}
{$800\text{ g}$}
{\True $850\text{ g}$}
{$900\text{ g}$}
\loigiai{
Đổi $1\text{ kg } 200\text{ g} = 1200\text{ g}$. Số táo còn lại: $1200 - 350 = 850\text{ (g)}$.
\textbf{Đáp án của câu hỏi là C}
}
\end{ex}

\begin{ex} %[3D3V34]
Để pha nước tắm cho bé sơ sinh vào mùa đông, nhiệt độ nước thích hợp nhất là khoảng
\choice
{\True $37^\circ\text{C} \text{ đến } 38^\circ\text{C}$}
{$20^\circ\text{C}$}
{$60^\circ\text{C}$}
{$100^\circ\text{C}$}
\loigiai{
Nhiệt độ nước tắm trẻ sơ sinh ấm ngang thân nhiệt, khoảng $37^\circ\text{C} - 38^\circ\text{C}$.
\textbf{Đáp án của câu hỏi là A}
}
\end{ex}

\begin{ex} %[3D3V34]
Dùng ca $500\text{ ml}$ múc nước đổ vào xô $10\ell$. Hỏi cần múc bao nhiêu ca đầy để đổ đầy xô?
\choice
{$10$ ca}
{$15$ ca}
{$25$ ca}
{\True $20$ ca}
\loigiai{
Đổi $10\ell = 10000\text{ ml}$. Số ca cần là: $10000 : 500 = 20$ ca.
\textbf{Đáp án của câu hỏi là D}
}
\end{ex}

\begin{ex} %[3D3V34]
Một chiếc bánh kem nặng $1\text{ kg}$. Người ta cắt thành 8 miếng đều nhau. Mỗi miếng bánh kem nặng là
\choice
{$150\text{ g}$}
{\True $125\text{ g}$}
{$100\text{ g}$}
{$200\text{ g}$}
\loigiai{
Đổi $1\text{ kg} = 1000\text{ g}$. Mỗi miếng bánh: $1000 : 8 = 125\text{ (g)}$.
\textbf{Đáp án của câu hỏi là B}
}
\end{ex}

\begin{ex} %[3D3V34]
Chiều cao của một bạn học sinh lớp 3 thường vào khoảng
\choice
{$130\text{ mm}$}
{$130\text{ m}$}
{\True $130\text{ cm}$}
{$13\text{ m}$}
\loigiai{
Học sinh lớp 3 cao khoảng $125\text{ cm} - 135\text{ cm}$.
\textbf{Đáp án của câu hỏi là C}
}
\end{ex}

\begin{ex} %[3D3V34]
Một gói bánh quy có 10 chiếc bánh, cả gói nặng $250\text{ g}$. Vỏ gói nặng $10\text{ g}$. Mỗi chiếc bánh quy nặng là
\choice
{\True $24\text{ g}$}
{$25\text{ g}$}
{$20\text{ g}$}
{$26\text{ g}$}
\loigiai{
Khối lượng của 10 chiếc bánh: $250 - 10 = 240\text{ g}$. Mỗi chiếc bánh: $240 : 10 = 24\text{ (g)}$.
\textbf{Đáp án của câu hỏi là A}
}
\end{ex}

\textbf{Phần 2. Trắc nghiệm đúng sai}
\setcounterpageref{ex}{0}

\begin{ex} %[3D3V34]
Xét tính đúng sai của các ước lượng thực tế:
\choiceTF
{\True Một gói muối ăn gia đình thường có khối lượng khoảng $500\text{ g}$ hoặc $1\text{ kg}$.}
{Một ly nước uống bình thường chứa khoảng $5\ell$ nước.}
{\True Bề dày của một cuốn vở học sinh khoảng $3\text{ mm}$ đến $5\text{ mm}$.}
{Nhiệt độ ngoài trời khi trời nắng gắt mùa hè ở miền Bắc khoảng $80^\circ\text{C}$.}
\loigiai{
\begin{itemize}
    \item Ý a) Đúng: Thực tế sản phẩm thị trường.
    \item Ý b) Sai: Sai vì ly nước chỉ chứa khoảng $200\text{ ml} - 300\text{ ml}$.
    \item Ý c) Đúng: Thực tế đo lường.
    \item Ý d) Sai: Sai vì nhiệt độ khoảng $38^\circ\text{C} - 42^\circ\text{C}$, $80^\circ\text{C}$ là gần sôi.
\end{itemize}
\begin{itemize}
    \item a) \textbf{Đúng}.
    \item b) \textbf{Sai}.
    \item c) \textbf{Đúng}.
    \item d) \textbf{Sai}.
\end{itemize}
}
\end{ex}

\begin{ex} %[3D3V34]
Bạn An thực hiện đo các đồ vật trong lớp học:
\choiceTF
{\True An dùng thước dây để đo chu vi thân cây bàng ở sân trường.}
{\True An dùng ê-ke để kiểm tra góc bảng lớp có vuông hay không.}
{An dùng nhiệt kế để đo khối lượng cặp sách.}
{\True An dùng ca có vạch chia để đo lượng nước trong bình.}
\loigiai{
\begin{itemize}
    \item Ý a) Đúng: Thước dây mềm uốn cong đo chu vi cây.
    \item Ý b) Đúng: Ứng dụng góc vuông của ê-ke.
    \item Ý c) Sai: Sai vì nhiệt kế đo nhiệt độ, không đo khối lượng.
    \item Ý d) Đúng: Ca chia vạch đo dung tích chất lỏng.
\end{itemize}
\begin{itemize}
    \item a) \textbf{Đúng}.
    \item b) \textbf{Đúng}.
    \item c) \textbf{Sai}.
    \item d) \textbf{Đúng}.
\end{itemize}
}
\end{ex}

\textbf{Phần 3. Trả lời ngắn}
\setcounterpageref{ex}{0}

\begin{ex} %[3D3V34]
Một cục tẩy dài $3\text{ cm } 5\text{ mm}$. Hỏi độ dài của cục tẩy là bao nhiêu mi-li-mét?
\shortans{35}
\loigiai{
\textbf{Đáp số:} $35$\\
$35\text{ mm}$.
}
\end{ex}

\begin{ex} %[3D3V34]
Một quả bưởi nặng $1\text{ kg } 400\text{ g}$. Hỏi quả bưởi đó nặng bao nhiêu gam?
\shortans{1400}
\loigiai{
\textbf{Đáp số:} $1400$\\
$1400\text{ g}$.
}
\end{ex}

\begin{ex} %[3D3V34]
Một chai nước ngọt có dung tích $1\ell\text{ }500\text{ ml}$. Hỏi chai nước ngọt đó chứa bao nhiêu mi-li-lít?
\shortans{1500}
\loigiai{
\textbf{Đáp số:} $1500$\\
$1500\text{ ml}$.
}
\end{ex}

\begin{ex} %[3D3V34]
Buổi chiều nhiệt độ đo được là $34^\circ\text{C}$, buổi tối nhiệt độ giảm xuống còn $26^\circ\text{C}$. Hỏi nhiệt độ đã giảm bao nhiêu độ C?
\shortans{8}
\loigiai{
\textbf{Đáp số:} $8$\\
$34 - 26 = 8^\circ\text{C}$.
}
\end{ex}

\begin{ex} %[3D3V34]
Một bao gạo $50\text{ kg}$ chia đều vào các túi nhỏ, mỗi túi chứa $5\text{ kg}$. Hỏi chia được bao nhiêu túi gạo?
\shortans{10}
\loigiai{
\textbf{Đáp số:} $10$\\
$50 : 5 = 10$ túi.
}
\end{ex}

\begin{ex} %[3D3V34]
Dùng một thìa đong $10\text{ ml}$ để đong đầy một bình $150\text{ ml}$. Hỏi cần đong bao nhiêu thìa?
\shortans{15}
\loigiai{
\textbf{Đáp số:} $15$\\
$150 : 10 = 15$ thìa.
}
\end{ex}
